\SourceLeaf{43}devoted to it. The law there is expressed with less than full generality. I thought it worth specifying it in more concise notation and deriving it from its true source by a rigorous proof.

First I shall briefly set out some preliminary remarks on a new sign, very useful for expressing the general terms of series. In series, particularly the more complicated ones, it often happens that odd and even terms appear to follow different laws and are then commonly split into two separate series. Such a division cannot be made without lengthening the expression, and it destroys, so to speak, the unity of expression that alone makes an indefinite series tractable. Unless some remedy is provided, it will confine the use of series within narrow limits. The more general method I now propose seems to arise naturally from the method of incorporating into the general term itself the simplest possible difference between even and odd terms of a series, namely that one set is positive and the other negative. Let the general term of a series be $m^{\text{tus}} = Q$. If you want it to be positive when $m$ is even and negative when $m$ is odd, I write it with the factor $(-1)^m$ prefixed. If you prefer the signs to proceed in the opposite order, I prefix $(-1)^{m+1}$\HistoricalNote{*)}{*)}{The distinguished\ \textsc{Kaestner} was kind enough not to withhold his approval of this sign when I proposed it to him. Relying on his authority, I have already used it several times in analytical investigations.}. Likewise, let the general expression for term $m^{\text{ti}}$, when $m$ is even, be $= P$, and when it is odd, be $= Q$, so that $P$ follows a law quite different from that of $Q$. I say that term $m^{\text{tum}}$ will generally be $\frac{1-(-1)^m}{2}Q+\frac{1-(-1)^{m+1}}{2}P$. For in this expression, if $m$ is even, $(-1)^m=1$ gives $\frac{1-(-1)^m}{2}=0$, and $(-1)^{m+1}=-1$ gives $\frac{1-(-1)^{m+1}}{2}=1$, so the first part vanishes and the other becomes $=P$, as required. If $m$ is odd, then $(-1)^{m+1}=1$ and hence $\frac{1-(-1)^{m+1}}{2}=0$, while $(-1)^m=-1$ and therefore $\frac{1-(-1)^m}{2}=1$. The latter part vanishes and the former remains $=Q$, as assumed. To give just one example that makes the usefulness of the notation clear, I shall assume what will be shown below, that an arbitrary power of the cosine of a given angle, e.\,g.\ $\cos\varphi^n$, can always be expanded in a series of cosines of multiples of the same angle. The law is that, when $n$ is even, the series proceeds by even multiples, so that its term $r^{\text{tus}}$, with the coefficient omitted, is $\cos 2r\varphi$. When $n$ is odd, it proceeds by odd multiples and its term is then $r^{\text{tus}} = \cos(2r+1)\varphi$. The general expression can accordingly be given as follows. In the series into which $\cos\varphi^n$ is%
